\documentclass[10pt,a4paper]{article}

\usepackage[utf8]{inputenc}
\usepackage[T1]{fontenc}
\usepackage[a4paper,top=0.8cm,bottom=0.9cm,left=0.7cm,right=0.7cm]{geometry}
\usepackage{amsmath,amssymb}
\usepackage{color}
\usepackage{multicol}

% ---------- Farben (dunkles Rot, druck- & cremepapier-tauglich) ----------
\definecolor{head}{rgb}{0.50,0.04,0.08}   % sehr dunkles Rot (Überschriften)
\definecolor{sub}{rgb}{0.62,0.07,0.10}    % dunkles Rot (Unterüberschriften)
\definecolor{grau}{rgb}{0.36,0.36,0.36}   % dunkleres Grau (Notizen)
\definecolor{vgcol}{rgb}{0.62,0.07,0.10}  % Marker = dunkelrot (statt Orange)

% ---------- Layout ----------
\setlength{\parindent}{0pt}
\setlength{\parskip}{2.8pt}
\setlength{\columnsep}{10pt}
\setlength{\columnseprule}{0.3pt}
\setlength{\abovedisplayskip}{3.5pt}
\setlength{\belowdisplayskip}{3.5pt}
\setlength{\abovedisplayshortskip}{1.5pt}
\setlength{\belowdisplayshortskip}{1.5pt}
\renewcommand{\baselinestretch}{1.18}
\pagestyle{empty}

% ---------- Überschriften (Referenz-Stil: Text + Linie) ----------
\newcommand{\Sec}[1]{\par\addvspace{5pt}%
  {\color{head}\large\bfseries #1}\par\vspace{-1.5pt}%
  {\color{head}\rule{\columnwidth}{0.9pt}}\par\vspace{1.5pt}}
\newcommand{\Sub}[1]{\par\addvspace{3pt}{\color{sub}\bfseries #1}\nobreak\par\vspace{0.5pt}}
% inline-Rezept (grau, nach der Formel)
\newcommand{\rec}[1]{{\color{vgcol}\footnotesize$\blacktriangleright$}\,{\color{grau}\itshape #1}}
\newcommand{\hint}[1]{{\color{grau}\itshape #1}}

\newenvironment{tl}{\begin{list}{\textbullet}{%
  \setlength{\leftmargin}{1.0em}\setlength{\labelsep}{0.38em}%
  \setlength{\itemsep}{2.5pt}\setlength{\parsep}{0pt}\setlength{\topsep}{2pt}%
  \setlength{\partopsep}{0pt}}}{\end{list}}

\begin{document}
\normalsize

{\centering {\color{head}\large\bfseries Analysis II — Open-Book-Formelsammlung}\;
{\color{grau}\footnotesize ($\blacktriangleright$ = Vorgehen)}\par}
\vspace{2pt}

\begin{multicols}{3}

% ===================== Ableitungen =====================
\Sec{Ableitungen}
\begin{tl}
\item $(a\,x^{b})' = a\,b\,x^{b-1}$, \;\;$(a\ln x)'=\tfrac{a}{x}$, \;\;$(a\,e^{x})'=a\,e^{x}$
\item Produkt $(fg)'=f'g+fg'$; \;Quotient $\big(\tfrac fg\big)'=\tfrac{f'g-fg'}{g^2}$
\item Kette $\big(g(h(x))\big)'=h'(x)\,g'(h(x))$
\item $(\sin)'=\cos$, \;$(\cos)'=-\sin$, \;$(\tan)'=\tfrac{1}{\cos^2 x}=1+\tan^2 x$
\item $(\arcsin x)'=\tfrac{1}{\sqrt{1-x^2}}$, \;$(\arccos x)'=\tfrac{-1}{\sqrt{1-x^2}}$
\item $(\arctan x)'=\tfrac{1}{1+x^2}$
\item $(\sinh x)'=\cosh x$, \;$(\cosh x)'=\sinh x$, \;$(\tanh x)'=\tfrac{1}{\cosh^2 x}$
\item $(A\sin\omega t)'=A\omega\cos\omega t$, \;$(A\cos\omega t)'=-A\omega\sin\omega t$
\end{tl}

% ===================== Vektoranalysis =====================
\Sec{Vektoranalysis}
\Sub{Totales Differential / Gradient}
$df=\partial_x f\,dx+\partial_y f\,dy$, \;\;$|\vec v|=\sqrt{v_1^2+\dots+v_n^2}$\par
$\nabla f=(\partial_{x_1}f,\dots,\partial_{x_n}f)^{\!\top}$
\Sub{Richtungsableitung}
$D_{\vec v}f=\nabla f\cdot\dfrac{\vec v}{|\vec v|}$ \quad
\rec{$\nabla f$ $\to$ Punkt $\to$ $\vec v$ normieren $\to$ Skalarprodukt}
\Sub{Divergenz \& Rotation}
$\operatorname{div}\vec f=\partial_x f_1+\partial_y f_2(+\partial_z f_3)$\par
$\operatorname{rot}_{2D}\vec f=\partial_x f_2-\partial_y f_1$ \hint{(Skalar)}
\[ \operatorname{rot}_{3D}\vec f=\begin{pmatrix}\partial_y f_3-\partial_z f_2\\ \partial_z f_1-\partial_x f_3\\ \partial_x f_2-\partial_y f_1\end{pmatrix} \]
\rec{Quellenfrei $\operatorname{div}=0$; wirbelfrei $\operatorname{rot}=0$; Wurzel $\Rightarrow\pm$}
\Sub{Weitere}
$\vec a\times\vec b=(a_2b_3-a_3b_2,\,a_3b_1-a_1b_3,\,a_1b_2-a_2b_1)^{\!\top}$\par
$\Delta f=\operatorname{div}(\operatorname{grad}f)=\sum_i\partial_{x_i}^2 f$\par
$\operatorname{rot}(\operatorname{grad}f)=\vec0$, \;\;$\operatorname{div}(\operatorname{rot}\vec f)=0$\par
Extremstelle: $\nabla f=\vec0$ \rec{System mit Einsetzen lösen}

% ===================== DGL =====================
\Sec{Differentialgleichungen}
\Sub{Linear homogen \;$y'=a(x)\,y$}
$y=c\,e^{A(x)}$, \;$A=\int a\,dx$. \;Zerfall $y'=\alpha y\Rightarrow y_0e^{\alpha t}$, \;$t_{1/2}=\tfrac{\ln2}{-\alpha}$
\Sub{Linear inhomogen \;$y'=a y+b$}
$y=e^{A}\big(y_0+\int_{x_0}^{x} b(u)\,e^{-A(u)}du\big)$\par
\rec{$a,b$ $\to$ $A=\int a$ $\to$ $e^{A},e^{-A}$ kürzen $\to$ Integral m. Grenzen $\to$ einsetzen}
\Sub{Trennung \;$y'=g(x)\,h(y)$}
$\int\tfrac{1}{h(y)}dy=\int g(x)\,dx$ \;\rec{nach $y$ auflösen $\to$ $c$ aus Startwert}

% ===================== Integralrechnung =====================
\Sec{Integralrechnung}
$\int_a^b f\,dx=F(b)-F(a)$
\begin{tl}
\item $\int a x^b dx=\tfrac{a}{b+1}x^{b+1}\,(b\!\neq\!-1)$, \;$\int\tfrac1x dx=\ln|x|$
\item $\int e^x dx=e^x$, \;$\int f'e^{f}dx=e^{f}$, \;$\int\ln x\,dx=x\ln x-x$
\item $\int\sin x\,dx=-\cos x$, \;$\int\cos x\,dx=\sin x$
\end{tl}
\Sub{Partielle Integration}
$\int uv'\,dx=uv-\int u'v\,dx$ \;\rec{$u=$ Potenz/$\ln$, $v'=$ Rest; Produkt bleibt $\to$ nochmal}
\Sub{Substitution}
$\int f(g)\,g'\,dx=\int f(u)\,du$ \;\rec{$u=$ innere Fkt., $du=u'dx$, kürzen, integrieren, zurück}
\Sub{Doppelintegral}
$\int_{a_y}^{b_y}\!\int_{a_x}^{b_x} f\,dx\,dy$ \;\rec{innen nach $x$ ($y$ konst., $\int c\,dx=cx$), dann nach $y$}

% ===================== Trigonometrie =====================
\Sec{Trigonometrie}
$\sin t,\cos t:\mathbb R\to[-1,1]$, Periode $2\pi$, \;$\cos t=\sin(t+\tfrac\pi2)$\par
$\sin^2 t+\cos^2 t=1$, \;$\tan t=\tfrac{\sin t}{\cos t}$, \;$1+\tan^2 t=\tfrac{1}{\cos^2 t}$\par
$\cosh^2 t-\sinh^2 t=1$, \;$\tanh t=\tfrac{\sinh t}{\cosh t}$\par
$\sec=\tfrac1{\cos}$, \;$\csc=\tfrac1{\sin}$, \;$\cot=\tfrac1{\tan}$\par
Welle: $\lambda=\tfrac1f=\tfrac{2\pi}{\omega}$, \;$\omega=2\pi f$\par
$\sin\!:0,1,0$ \;$\cos\!:1,0,-1$ \hint{$(0,\tfrac\pi2,\pi)$}\par
\rec{Trig-Integral: $\tan\!\to\!\tfrac{\sin}{\cos}$, ausklammern, $\sin^2+\cos^2=1$, integrieren}

% ===================== Komplexe Zahlen =====================
\Sec{Komplexe Zahlen}
$i^2=-1$, \;$z=a+ib$, \;$\operatorname{Re}z=a$, \;$\operatorname{Im}z=b$, \;$\overline z=a-ib$\par
$|z|=\sqrt{a^2+b^2}$ \hint{(beide quadrieren!)}\par
$z_1z_2=(a_1a_2-b_1b_2)+i(a_1b_2+a_2b_1)$\par
$z^2=(a^2-b^2)+2ab\,i$\par
$\dfrac{z_1}{z_2}=\dfrac{(a_1a_2+b_1b_2)+i(a_2b_1-a_1b_2)}{a_2^2+b_2^2}$ \rec{mit $\overline z$ erweitern}\par
$z=r(\cos\varphi+i\sin\varphi)=r e^{i\varphi}$, \;$r=|z|$, \;$\varphi=\arctan\tfrac ba$\par
$e^{x+iy}=e^{x}(\cos y+i\sin y)$, \;$z^x=r^x e^{i\varphi x}$, \;$e^{i\pi}+1=0$

% ===================== Fourier =====================
\Sec{Fourier}
$A(\cos\omega t+i\sin\omega t)=A e^{i\omega t}$\par
$F(\omega)=\int_{-\infty}^{\infty} f(t)e^{-i\omega t}dt$, \;\;
$f(t)=\tfrac{1}{2\pi}\int_{-\infty}^{\infty}F(\omega)e^{i\omega t}d\omega$

\end{multicols}
\end{document}
